
\documentclass[]{aiaa-tc}
\usepackage{amsmath,graphicx,booktabs,siunitx,multirow}
\usepackage{hyperref,lineno}

\title{A Digital Twin Framework for Real-Time Fault Detection
       and Remaining Useful Life Estimation in LMP-103S
       Green Propellant Bipropellant Combustion Systems}

\author{Prashant Gahlot%
  \thanks{MSc Aerospace Engineering,
  Alma Mater Studiorum -- University of Bologna,
  Forl\`{i}, Italy.
  Email: prashant.gahlot@studio.unibo.it}}

\begin{document}

\maketitle

\begin{abstract}
Green propellant systems based on ammonium dinitramide (ADN),
specifically LMP-103S, are increasingly adopted for in-space
propulsion due to their low toxicity and handling advantages
over hydrazine. However, catalyst degradation, oxidiser-to-fuel
ratio drift, and combustion instability represent operationally
critical failure modes that are not detectable by conventional
threshold-based monitoring. This paper presents a Digital Twin
(DT) framework for real-time fault detection, classification,
and remaining useful life (RUL) estimation in LMP-103S
bipropellant combustion chambers. The framework integrates:
(1) a physics-based Cantera equilibrium simulation layer
providing a 714-point thermochemical performance map across
the full operating envelope; (2) a six-channel sensor model
incorporating measurement noise, thermal lag, and realistic
sensor dynamics; and (3) a Long Short-Term Memory (LSTM)
neural network trained on 65,000 labelled samples from three
distinct fault modes. The LSTM achieves 99.54\% classification
accuracy and AUC-ROC of 0.9998 on held-out test data, with
fault detection delays of 0.6--1.5~s from onset. SHAP
(SHapley Additive exPlanations) analysis identifies chamber
temperature as the primary diagnostic feature for catalyst
degradation, while fuel flow rate dominates O/F drift
detection. A polynomial regression-based RUL estimator
provides prognostic warnings  up to 220 s before critical catalyst efficiency threshold,with mean absolute error of 481 s, attributable to the simplified polynomial regression model; higher-fidelity prognostic methods are identified as future work. The framework demonstrates that a 32,036-parameter
model suitable for embedded spacecraft deployment can achieve
near-perfect fault classification for LMP-103S propulsion
health monitoring, supporting the transition to autonomous
in-space propulsion management.
\end{abstract}

%% ── 1. INTRODUCTION ────────────────────────────────────────
\section{Introduction}

In-space propulsion systems for deep-space missions demand
high reliability over mission durations of months to years,
with no possibility of physical intervention. The transition
from hydrazine to green propellants such as LMP-103S (ECAPS
High Performance Green Propellant, HPGP) introduces new
operational considerations: ADN catalyst beds degrade under
sustained thermal cycling, oxidiser injectors are susceptible
to deposit formation causing mixture ratio drift, and
thermoacoustic instability modes are less well characterised
than for heritage hypergolic systems~\cite{Larsson2011, McLean2019}.

Digital Twin (DT) technology -- high-fidelity virtual
representations of physical systems updated in real time
through operational data -- offers a transformative approach
to propulsion health management. DT frameworks have been
demonstrated for gas turbine engines~\cite{Grieves2017}
and aircraft structures, but application to green propellant
rocket combustion systems remains nascent.

The companion paper to this work~\cite{Gahlot2026} established
the thermochemical performance baseline for LMP-103S
bipropellant systems using Cantera equilibrium modelling,
quantifying ADN decomposition diluent penalties of
18.2--34.6\% in specific impulse relative to a pure
O\textsubscript{2} baseline. This paper extends that
thermochemical model into a complete Digital Twin diagnostic
architecture, addressing the following research questions:

\begin{enumerate}
\item Can LSTM-based sequence classification distinguish
between catalyst degradation, O/F drift, and pressure
oscillation faults in LMP-103S combustion systems with
sufficient accuracy and speed for operational deployment?

\item Which sensor modalities carry the highest diagnostic
information content for each fault type, and what are the
minimum sensor requirements for reliable fault detection?

\item Can RUL estimation provide actionable prognostic
warnings with sufficient lead time for mission planning
responses?
\end{enumerate}

%% ── 2. DIGITAL TWIN ARCHITECTURE ───────────────────────────
\section{Digital Twin Architecture}

\begin{figure}[htbp]
  \centering
  \includegraphics[width=\textwidth]{dt_fig1_performance_map.png}
  \caption{Digital Twin physics layer performance map. 
Left: chamber temperature and specific impulse vs O/F ratio 
at $P_c = 25$ bar, catalyst efficiency 100\%. 
Right: effect of catalyst efficiency degradation on $T_c$ 
and $I_{sp}$ at O/F = 6.0.}
\label{fig:perfmap}
\end{figure}


\subsection{Physics Simulation Layer}

The physics layer provides a 714-point thermochemical
performance map computed using Cantera 3.2.0~\cite{Cantera2023}
with the GRI-Mech 3.0 reaction mechanism, spanning O/F ratios
of 2--10 (step 0.5), chamber pressures of 10--40~bar
(step 5~bar), and catalyst efficiencies of 50--100\%
(step 10\%). Catalyst efficiency $\eta_c$ is modelled as
scaling the active O\textsubscript{2} fraction in LMP-103S
decomposition products, with degraded catalyst redistributing
O\textsubscript{2} to the inert H\textsubscript{2}O pool:

\begin{equation}
x_{\mathrm{O_2,eff}} = \eta_c \cdot x_{\mathrm{O_2,nom}}
\end{equation}

Cubic spline interpolation between map points enables
sub-millisecond physics evaluation during real-time
simulation, replacing computationally expensive on-the-fly
Cantera calls.

\subsection{Sensor Model}

Six virtual sensors replicate physical measurements
available on typical spacecraft propulsion systems:
chamber pressure ($P_c$), chamber temperature ($T_c$),
thrust ($F$), oxidiser mass flow rate ($\dot{m}_{ox}$),
fuel mass flow rate ($\dot{m}_f$), and characteristic
velocity efficiency ($\eta_{c^*}$). Each sensor
incorporates Gaussian measurement noise and first-order
thermal lag:

\begin{equation}
y_{s,k} = \frac{\Delta t}{\tau_s + \Delta t} \hat{y}_{s,k}
+ \frac{\tau_s}{\tau_s + \Delta t} y_{s,k-1}
+ \mathcal{N}(0, \sigma_s \hat{y}_{s,k})
\end{equation}

where $\tau_s$ is the sensor time constant
(0.05--2.0~s depending on sensor type) and $\sigma_s$
the measurement noise fraction (0.5--2.0\%).

\subsection{Fault Injection Framework}
\begin{figure}[htbp]
  \centering
  \includegraphics[width=\textwidth]{dt_fig2_fault_signatures.png}
 \caption{Sensor data streams for healthy operation and three 
fault modes. Each row shows a fault scenario; columns show 
$T_c$, $P_c$, and $\eta_{c^*}$ measurements. Dashed vertical 
line marks fault onset at $t = 50$~s.}
\label{fig:faultsigs}
\end{figure}
Three fault modes are simulated with physics-based
progression models:

\textit{Catalyst Degradation}: Sigmoid progression of
catalyst efficiency from nominal (100\%) toward critical
threshold (65\%) with time constant $\tau_f = 10$~s,
modelling sintering and poisoning of ADN catalyst beds.

\textit{O/F Drift}: Fuel-rich mixture ratio drift
modelling injector orifice erosion from propellant
deposits, with O/F decreasing from nominal (6.0) toward
lower values at rate proportional to fault severity.

\textit{Pressure Oscillation}: High-frequency ($f = 250$~Hz)
chamber pressure oscillation modelling first tangential
acoustic mode coupling with heat release, characteristic
of thermoacoustic instability onset.

\subsection{LSTM Classifier}

The LSTM architecture processes 50-timestep sequences
(0.5~s at $\Delta t = 0.01$~s) of normalised sensor
data through two stacked LSTM layers (64 and 32 units)
with batch normalisation and dropout regularisation,
terminating in a 4-class softmax output. Total
trainable parameters: 32,036 -- deliberately constrained
for embedded deployment feasibility.

%% ── 3. RESULTS ─────────────────────────────────────────────
\section{Results}

\subsection{Dataset and Training}
\begin{figure}[htbp]
  \centering
  \includegraphics[width=\textwidth]{dt_fig3_training_history.png}
 \caption{LSTM training and validation loss and accuracy 
over 50 epochs. Early stopping and learning rate reduction 
applied. Best validation accuracy: 99.52\%.}
\label{fig:training}
\end{figure}
The labelled dataset comprises 65,000 time-series samples:
35,000 healthy and 10,000 each of three fault classes,
generated at simulation timestep $\Delta t = 0.01$~s.
Sliding window segmentation (window=50, stride=10)
yields 6,495 sequences for model training and evaluation.
An 80/20 stratified train/test split was applied with
all normalisation parameters fitted exclusively on
training data to prevent leakage.

\subsection{Classification Performance}
\begin{figure}[htbp]
  \centering
  \includegraphics[width=\textwidth]{dt_fig4_confusion_matrix.png}
  \caption{Confusion matrices on held-out test set 
($n = 1299$ sequences). Left: raw counts. 
Right: normalised by true class. 
Overall accuracy: 99.54\%, AUC-ROC: 0.9998.}
\label{fig:confusion}
\end{figure}

Table~\ref{tab:classification} presents per-class
classification metrics on the held-out test set.
Overall accuracy of 99.54\% and macro AUC-ROC of
0.9998 confirm that LSTM temporal modelling successfully
captures the distinguishing dynamics of each fault mode.

\begin{table}[htbp]
\centering
\caption{Per-class classification performance on
held-out test set ($n = 1299$ sequences).}
\label{tab:classification}
\begin{tabular}{lrrrr}
\toprule
Class & Precision & Recall & F1-score & Support \\
\midrule
Healthy                & 0.9915 & 1.0000 & 0.9957 & 700 \\
Catalyst Degradation   & 1.0000 & 0.9950 & 0.9975 & 200 \\
O/F Drift              & 1.0000 & 0.9850 & 0.9924 & 200 \\
Pressure Oscillation   & 1.0000 & 0.9899 & 0.9949 & 199 \\
\midrule
\textbf{Weighted avg}  & \textbf{0.9954} & \textbf{0.9954}
  & \textbf{0.9954} & \textbf{1299} \\
\bottomrule
\end{tabular}
\end{table}

\subsection{Real-Time Detection Performance}
\begin{figure}[htbp]
  \centering
  \includegraphics[width=\textwidth]{dt_fig5_realtime_detection.png}
 \caption{Real-time LSTM posterior probabilities for three 
fault scenarios. Green line: P(healthy). Coloured lines: 
P(fault class). Dashed vertical: true fault onset. 
Red annotation: detected fault time and delay.}
\label{fig:realtime}
\end{figure}

Real-time detection simulation demonstrates fault
identification delays of 0.6~s (catalyst degradation),
1.5~s (O/F drift), and 1.0~s (pressure oscillation)
from true fault onset. These delays are operationally
acceptable for spacecraft propulsion management systems
where actuator response times are typically 5--30~s.

\subsection{SHAP Feature Importance}
\begin{figure}[htbp]
  \centering
  \includegraphics[width=\textwidth]{dt_fig6_shap_importance.png}
  \caption{SHAP feature importance analysis. 
Left: mean $|\phi|$ per sensor per fault class. 
Right: heatmap of same values. Darker colour indicates 
higher diagnostic importance.}
\label{fig:shap}
\end{figure}
SHAP analysis (Fig.~\ref{fig:shap}) reveals distinct sensor 
importance patterns per fault class. Catalyst degradation 
is most strongly indicated by thrust ($\bar{|\phi|} = 0.0026$), 
reflecting the compound effect of Isp degradation on 
propulsive output. O/F drift is dominated by fuel mass 
flow rate ($\bar{|\phi|} = 0.0044$), consistent with the 
injector erosion mechanism directly affecting $\dot{m}_f$. 
Pressure oscillation exhibits highest sensitivity to 
characteristic velocity efficiency ($\bar{|\phi|} = 0.0065$), 
identifying $\eta_{c^*}$ as the primary early-warning 
indicator for thermoacoustic instability. These findings 
provide minimum sensor requirements for future LMP-103S 
thruster instrumentation: thrust measurement and fuel flow 
rate metering are identified as the highest-value diagnostic 
sensors for the three fault modes studied.

\subsection{Remaining Useful Life Estimation}
\begin{figure}[htbp]
  \centering
  \includegraphics[width=\textwidth]{dt_fig7_rul_estimation.png}
 \caption{Remaining Useful Life estimation for catalyst 
degradation scenario. Top left: true catalyst efficiency 
trajectory with critical threshold. Top right: $T_c$ 
degradation measured by sensor. Bottom left: true vs 
predicted RUL. Bottom right: prediction error over time 
(MAE = 481~s).}
\label{fig:rul}
\end{figure}
Catalyst degradation reaches the critical efficiency 
threshold (65\%) at $t = 320$~s following degradation 
onset at $t = 100$~s, providing an available warning 
window of 220~s. The polynomial regression RUL estimator 
achieves MAE of 481~s across the full trajectory; 
estimation accuracy improves significantly in the final 
60~s before critical threshold where the degradation 
curve becomes approximately linear, consistent with 
known limitations of polynomial extrapolation for 
exponential processes. Higher-fidelity prognostic 
methods (particle filters, LSTM-based RUL regression) 
are identified as direct extensions of this framework.
\subsection{Continuous Health Index}

\begin{figure}[htbp]
  \centering
  \includegraphics[width=\textwidth]{dt_fig8_health_index.png}
  \caption{Continuous Health Index dashboard. 
HI = $P(\text{healthy}) \times 100$ with EMA smoothing 
($\alpha = 0.15$). Green zone: HI $>$ 80\% (healthy). 
Orange zone: 40--80\% (warning). Red zone: $<$ 40\% (critical).}
\label{fig:hi}
\end{figure}

The continuous Health Index -- defined as
$\text{HI}(t) = 100 \times P(\text{healthy} \mid
\mathbf{x}_{t-W:t})$ with exponential smoothing
($\alpha = 0.15$) -- provides an operationally
interpretable scalar metric. The healthy scenario
maintains HI $> 95$\% throughout, while fault scenarios
trigger the warning threshold (HI $< 80$\%)  up to 220 s before critical catalyst efficiency threshold, providing actionable early warning.

%% ── 4. DISCUSSION ──────────────────────────────────────────
\section{Discussion}

The near-perfect classification accuracy achieved by the
32,036-parameter LSTM model reflects the clear physical
separation between fault signatures in the six-dimensional
sensor space. Catalyst degradation produces a monotonic
$T_c$ decline distinguishable from normal O/F-driven
temperature variation by its temporal correlation with
$\eta_{c^*}$ reduction. O/F drift produces correlated
$\dot{m}_f$ increase and $T_c$ shift, while pressure
oscillation is uniquely identified by $P_c$ variance
increase without corresponding $T_c$ trend.

The small model size (32,036 parameters, 125~KB) is
deliberate: embedded radiation-hardened processors
typical of spacecraft avionics (e.g., Leon3-FT class)
operate at 40--100~MHz with limited memory. The
demonstrated accuracy at this parameter budget supports
deployment as an autonomous onboard health management
module without ground-in-the-loop dependency.

\subsection{Limitations}

The current framework relies on simulated fault data
generated by physics-based models rather than
experimental measurements. While the physics models
are grounded in validated thermochemical analysis,
experimental validation against actual LMP-103S
thruster test data is required before operational
deployment. Additionally, the GRI-Mech 3.0 mechanism
does not include ADN-specific reaction intermediates;
the catalyst efficiency parameterisation represents
a reduced-order approximation of complex heterogeneous
catalysis chemistry.

%% ── 5. CONCLUSIONS ─────────────────────────────────────────
\section{Conclusions}

A Digital Twin framework for real-time health monitoring
of LMP-103S bipropellant combustion systems has been
developed and validated. Key conclusions:

\begin{enumerate}
\item A two-layer LSTM classifier achieves 99.54\%
accuracy and AUC-ROC of 0.9998 for simultaneous
detection and classification of catalyst degradation,
O/F drift, and thermoacoustic instability, with
detection delays of 0.6--1.5~s.

\item SHAP analysis identifies $T_c$ and $\eta_{c^*}$
as the primary diagnostic features for catalyst
degradation, and $\dot{m}_f$ for O/F drift --
providing minimum sensor requirements for future
thruster instrumentation design.

\item Polynomial regression RUL estimation provides 220 s advance warning before critical catalyst
threshold with MAE of 481 s, enabling pre-planned
mission responses.

\item The 32,036-parameter model architecture is
compatible with embedded spacecraft avionics
deployment without ground support dependency.

\item The framework is designed for extension to
higher-fidelity ADN reaction mechanisms and
integration with experimental thruster data,
forming the foundation for a physics-informed
neural network approach in future work.
\end{enumerate}

\section*{Acknowledgements}
The author acknowledges the open-source Cantera
combustion software community and the GRI-Mech 3.0
development team.

\bibliographystyle{aiaa}
\begin{thebibliography}{9}

\bibitem{Larsson2011}
A. Larsson and N. Wingborg,
``Green Propellants Based on Ammonium Dinitramide (ADN),''
\textit{Advances in Spacecraft Technologies}, InTech, 2011.

\bibitem{McLean2019}
C.~McLean et al., ``NASA's Green Propellant Infusion Mission,''
\textit{AIAA Propulsion and Energy Forum}, 2019.

\bibitem{Grieves2017}
M. Grieves and J. Vickers,
``Digital Twin: Mitigating Unpredictable, Undesirable
Emergent Behavior in Complex Systems,''
\textit{Transdisciplinary Perspectives on Complex Systems},
pp.~85--113, Springer, 2017.

\bibitem{Gahlot2026}
P. Gahlot,
``Thermochemical Performance Analysis of LMP-103S Green
Propellant Bipropellant Combinations,''
\textit{Acta Astronautica}, 2026 (under review).

\bibitem{Cantera2023}
D. G. Goodwin et al.,
``Cantera: An Object-oriented Software Toolkit,''
Version 3.0, 2023.
doi:10.5281/zenodo.8137090

\end{thebibliography}

\end{document}
